\section{An explicit CM cover killing a finite dihedral type}
\label{sec:cubic-dihedral-cm-cover}

\begin{proposition}[AC-G73: a controlled totally real ray-class cover]
\label{AC-G73}
Let $p>3$ and $q$ be odd primes with $q\equiv-1\pmod p$.
Let $\theta:\mathbb F_{q^2}^{\times}\to\overline{\mathbb Q}_p^{\times}$
be a character of exact order $p^r>1$.
There is a totally real Galois extension $L/\mathbb Q$ with
\[
 \operatorname{Gal}(L/\mathbb Q)\simeq C_{p^n}\rtimes C_2,
 \qquad n\ge r,
                                                        \tag{G73.1}
\]
where the involution acts by inversion, and a real quadratic
subfield $F$ such that $q$ is inert in $F$, $L/F$ is cyclic,
and its local inertia character at $q$ has kernel
$\ker\theta$ under the identification
$\mathcal O_F/q\mathcal O_F\simeq\mathbb F_{q^2}$.
The extension $L/F$ is unramified outside $q$ and at infinity.
One can arrange
\[
 D_F\ll q^{10},\qquad [L:\mathbb Q]\ll q^6\log(2q),
 \qquad \log D_L\ll[L:\mathbb Q]\log(2q),             \tag{G73.2}
\]
with absolute implied constants. No prime-location theorem
or RH is used in these bounds after $q$ is given.

If $K$ is an imaginary quadratic field in which $q$ splits,
then $B=KL$ is CM. Every absolutely irreducible
three-dimensional characteristic-zero representation of
$G_K$ remains absolutely irreducible on $G_B$.
If its inertia at both primes above $q$ factors through
$\theta\oplus\theta^{-1}\oplus1$ after the unramified
quadratic local extension, then its restriction to $G_B$
is unramified at those primes. This last assertion assumes
the actual Galois inertia statement; it does not supply
local--global compatibility or an automorphic seed.
\end{proposition}
\begin{proof}
\proofstep{A squarefree Pell construction with a small defining parameter.}
Choose $b_0\in\mathbb F_q^\times$ with $2b_0$ a nonsquare.
Choose an odd integer $b_*\in[1,2q]$ reducing to $b_0$.
For $1\le k\le T:=1024q^3$, put
\[
 b_k=b_*+2qk,\quad a_k=1+q^2b_k,\quad
 d_k=b_k(2+q^2b_k)=\frac{a_k^2-1}{q^2}.              \tag{G73.3}
\]
Both factors of $d_k$ are odd and their greatest common
 divisor is one. Neither is divisible by $q$.
For an odd prime $v\ne q$, the condition $v^2\mid d_k$
therefore holds in at most two residue classes modulo $v^2$.
Every such $v$ is at most
$\sqrt{2+q^2\max b_k}$, because it squares into one
of the two factors.
As $b_k\le2q(T+1)\le2050q^4$, this bound is at most
$46q^3$ for $q\ge3$.
The number of bad $k$ is consequently at most
\[
 2T\sum_{v\ge3}\frac1{v^2}+92q^3
 \le\frac89T+92q^3<T.                              \tag{G73.4}
\]
Here the sum was enlarged to all integers $v\ge3$, and
$\sum_{v\ge3}v^{-2}\le1/9+\int_3^\infty t^{-2}dt=4/9$.
Thus some $d=d_k$ is squarefree. Fix its $a=a_k$ and $b=b_k$.
It satisfies $d\equiv3\pmod4$, $d\ll q^{10}$, and
$d\equiv2b_0\pmod q$ is a nonsquare.
Therefore $F=\mathbb Q(\sqrt d)$ is real quadratic,
$\mathcal O_F=\mathbb Z[\sqrt d]$, $D_F=4d$, and $q$ is inert.

The unit
\[
                   \varepsilon=a+q\sqrt d           \tag{G73.5}
\]
has norm one and reduces to $1$ modulo $q$.
It is the fundamental positive unit.
Indeed norm-minus-one units do not exist when $d\equiv3\pmod4$,
as their norm equation modulo $4$ shows.
Every unit $u>1$ with positive conjugate has the form
$x+y\sqrt d$, $y\ge1$, and hence $u>2\sqrt d$.
If (G73.5) were a power $u^s$ with $s\ge2$, it would
exceed $4d$. But $\varepsilon<2a<4d$, because
\[
                  2d-a=q^2b(2b-1)+4b-1>0.
\]
The unit theorem now gives $\mathcal O_F^\times=\{\pm\varepsilon^m:m\in\mathbb Z\}$.
In particular every global unit has reduction in $\{\pm1\}$ at $q$.

\proofstep{Extend the full local character through the minus ray class group.}
Let $A$ be the $p$-primary part of the ray class group of
$F$ of finite modulus $q\mathcal O_F$, with no infinite
places in the modulus. The exact ray class sequence gives
an injection
\[
 C:=(\mathbb F_{q^2}^{\times})_p\lhook\joinrel\longrightarrow A,
                                                        \tag{G73.6}
\]
because the reductions of the global units have order at most two.
Let $\sigma$ be the nontrivial automorphism of $F/\mathbb Q$.
It acts on $C$ as $z\mapsto z^q$, which is inversion:
$p\nmid q-1$, so the whole $p$-part of $q^2-1$ divides $q+1$.
Since $p$ is odd, the projectors $(1\pm\sigma)/2$
split $A=A^+\oplus A^-$, and (G73.6) lies in $A^-$.

The restriction of $\theta$ to $C$ extends to a character
$\phi:A^-\to\mu_{p^\infty}$.
For an elementary construction, adjoin generators one at a time:
if the next generator $x$ has its least positive multiple
$p^s x$ in the existing subgroup, choose a $p^s$-th root
of the prescribed value at $p^s x$ and assign it to $x$.
This defines an extension, and the finite number of
steps leaves a finite $p$-power image.
Extend $\phi$ trivially on $A^+$ and on the prime-to-$p$
part of the ray class group. Its order is $p^n$ with $n\ge r$,
and $\phi^\sigma=\phi^{-1}$.
The global ray class correspondence gives its cyclic
extension $L/F$. Its finite conductor divides $q\mathcal O_F$,
and no infinite place ramifies; thus $L$ is totally real.
Local reciprocity identifies the inertia character with
$\theta$, up to inverse, which has the same kernel.

The kernel of $\phi$ is $\sigma$-stable, so $L/\mathbb Q$
is Galois and has quotient $C_2$ over $F$ acting by inversion
on $C_{p^n}$. For a lift $s$ of the involution,
$s^2\in C_{p^n}$ is fixed by that inversion, hence is one.
This proves (G73.1) directly, including the splitting.

\proofstep{Degree and discriminant cost.}
The exact ray class sequence gives
\[
 p^n\le |A|\le h_F(q+1)_p\le h_F(q+1).              \tag{G73.7}
\]
Minkowski's bound supplies an integral ideal of norm at
most $\sqrt{D_F}/2$ in each class of a real quadratic field.
There are at most $\tau(m)$ ideals of norm $m$ in a quadratic
field. Consequently
\[
 h_F\le\sum_{m\le\sqrt{D_F}/2}\tau(m)
       \ll\sqrt{D_F}\log(2D_F)\ll q^5\log(2q).
                                                        \tag{G73.8}
\]
This proves the degree bound in (G73.2).
Every character of $\operatorname{Gal}(L/F)$ has conductor
dividing $q\mathcal O_F$. The conductor--discriminant formula
therefore gives
\[
 D_L\le D_F^{p^n}q^{2(p^n-1)},                       \tag{G73.9}
\]
which gives the last assertion of (G73.2).
For fixed $K$, the same tower estimate gives
$\log D_B\ll_K[B:\mathbb Q]\log(2q)$.

\proofstep{The CM cover and restriction of an irreducible coefficient.}
The fields $K$ and $L$ are linearly disjoint because
$L$ is totally real and $K$ is imaginary quadratic.
Thus $B=KL$ is a totally imaginary quadratic extension
of the totally real field $L$, hence is CM.
There is a tower
\[
                   K\subset KF\subset KL
\]
of a quadratic cyclic extension followed by a cyclic
$p$-power extension. Absolute irreducibility of a
three-dimensional characteristic-zero representation
survives both steps. To see this without a residual
hypothesis, its restriction to a normal subgroup of
finite index is semisimple: conjugates of a simple
submodule span the original irreducible module.
For a cyclic quotient $C_m$, induction--restriction gives
\[
 \operatorname{End}_{G'}(V)
  =\bigoplus_{\eta\in\widehat{C_m}}
                    \operatorname{Hom}_{G}(V,V\otimes\eta).
                                                        \tag{G73.10}
\]
Every nonzero summand gives an isomorphism, so determinants
force $\eta^3=1$. For $m=2$ or a $p$-power, only the
trivial character is possible. The commutant of the
semisimple restriction is therefore one-dimensional,
so it remains absolutely irreducible.

At a prime of $K$ above $q$, adjoining $F$ supplies the
unramified quadratic local field. The further extension
$L/F$ has inertia kernel precisely $\ker\theta$.
The specified actual inertia action therefore becomes
trivial on $G_B$; this proves the last assertion.
\end{proof}

\paragraph{Sources and scope.}
The retained class-field inputs are Milne,
\emph{Class Field Theory}, v4.03 (2020), Chapter~V,
Theorem~1.7 and its complete proof (printed pp.~150--151),
Theorems~3.5--3.6 and Corollary~3.7 (p.~158), and
Theorem~3.27 (p.~165). The existence and conductor--discriminant
theorems are retained known results; their full upstream
proofs are not claimed to have been independently reproved.
Minkowski's exact form is the already audited Milne,
\emph{Algebraic Number Theory}, v3.08, Theorem~4.3, p.~70.
The squarefree construction, unit calculation, character
extension, quantitative consequences and restriction
argument above are the new proof.

This is a finite-cover construction for a prescribed
finite local character. It does not construct an actual
ordinary ramified automorphic point, prove positivity of
the G72 multiplicity difference, verify every away-prime
type, transport the old source frame to a new CM host,
or establish characteristic-$p$ curve support or H04.4.
Its finite-cover interpretation uses SGA~1, Expos\'e~V,
\S\S4 and~7 in the already audited scope.
